\documentclass[10pt,a4paper]{amsart}
\usepackage[margin=19mm]{geometry}
\usepackage{amsmath,amssymb,mathtools}
\usepackage[scr=rsfs,cal=euler]{mathalfa}
\usepackage{xfrac}
\usepackage[unicode,hidelinks,bookmarks=false]{hyperref}
\newcommand{\ie}{i.e.\ }
\newcommand{\cf}{cf.\ }
\newcommand{\resp}{respectively\ }
\newcommand{\Subsection}{\subsection}
\providecommand{\nobreakdash}{\nobreak\hbox{-}}
\setlength{\emergencystretch}{2em}
\multlinegap0pt
\usepackage{bm}
\usepackage{bbm}
\usepackage{dsfont}
\usepackage{xspace}
\usepackage{cancel}
\usepackage{mathtools}
\usepackage{amsthm}
\makeatletter
\@ifundefined{theorem}{\newtheorem{theorem}{Theorem}}{}
\@ifundefined{definition}{\newtheorem{definition}[theorem]{Definition}}{}
\@ifundefined{lemma}{\newtheorem{lemma}[theorem]{Lemma}}{}
\makeatother
\title[Inverse Hessian Curvature Flow in Minkowski Space II: The Di]{Inverse Hessian Curvature Flow in Minkowski Space II: The Dirichlet problem at infinity}
\author{Dake Li, Zhizhang Wang, Shiqi Yin}
\date{Source: arXiv:2608.20037v1 (CC BY 4.0)}
\begin{document}
\maketitle

\noindent\emph{Source and licence.} Inverse Hessian Curvature Flow in Minkowski Space II: The Dirichlet problem at infinity. Version arXiv:2608.20037v1 (CC BY 4.0), distributed under the Creative Commons Attribution 4.0 International licence, as recorded in the arXiv licence metadata for this exact version and shown on the abstract page https://arxiv.org/abs/2608.20037v1. Full rights notice: licenses/arxiv-2608.20037-CC-BY-4.0.txt.\par\smallskip

\noindent\textit{Editorial scope.} Counted unit: the Lemma whose source label is \texttt{Viscosity comparison principle} in the article (source0.tex lines 448--452, source label \texttt{Viscosity comparison principle}), delivered together with the article's own complete original proof of it (source0.tex lines 453--531, one \texttt{proof} environment of 79 lines). No mathematics is written, repaired or re-derived: statement and proof are the article's own text line for line. Required context retained uncounted: Comparison principle (lines 390--392); viscosity subsolution equation (lines 428--442); smooth function g (lines 1251--1257). Every definition, display and label of those lines is retained unchanged. Omitted from this package and not redistributed: 1--389, 393--427, 443--447, 532--1250, 1258--3011. Nothing inside a retained excerpt is omitted or altered: every retained body is delivered exactly as the article's lines are written, so no omission receipt of any kind accompanies this package. Citations: the retained excerpts carry no citation command at all, so no bibliography entry of the article is transcribed and no reference is redistributed. Reference adaptation: none was needed and none was made; every reference inside the delivered unit resolves inside it, so no unresolved reference or citation is produced. Printed slips kept literal: no printed word, symbol, hypothesis or number of the counted statement or of its proof was altered. Layout and class: the article's own class is \texttt{amsart} with packages including amsfonts, amssymb, amsmath, mathrsfs, amsthm, hyperref, enumitem, mathtools, soul; the wrapper is the lane's pinned \texttt{amsart} editorial wrapper at 10pt on A4, which restates the theorem-like environments the retained text needs and adds the article's own macro definitions that the retained text needs (none), with extra wrapper packages (bm, bbm, dsfont, xspace, cancel, mathtools). The delivered numbering is the unit's own. Assets: no retained span contains a figure environment, a graphics-inclusion command, a PDF-page inclusion, a table or a TikZ picture; no image file is copied into this package, the package declares no asset dependency and no image file is redistributed with it. Licence: version arXiv:2608.20037v1 of this article is distributed under the Creative Commons Attribution 4.0 International licence, as recorded in the arXiv licence metadata for this exact version and shown on the abstract page https://arxiv.org/abs/2608.20037v1. A full rights notice is delivered at licenses/arxiv-2608.20037-CC-BY-4.0.txt. Characters: every retained excerpt is pure ASCII, so the delivered engine, XeLaTeX with its default Latin Modern text font, reports no missing character.
\begin{lemma}[Comparison principle]\label{Comparison principle}
    Suppose $F=\left( \frac{\sigma_n}{\sigma_{n-k}} \right)^{\frac{1}{k}}$, then let $F(w)=F(w_{ij}-w\delta_{ij})$ by a slight abuse of notation. For $w,v \in \mathcal{S}$, if $F(w) + \frac{w}{\binom{n}{k}^{\frac{1}{k}}} \geqslant F(v) + \frac{v}{\binom{n}{k}^{\frac{1}{k}}}$ and $w \leqslant v$ on $\mathbb{S}^{n-1}$ in cone topology sense, we have $w \leqslant v$ in $\mathbb{H}^{n}$. 
\end{lemma}

\begin{definition}[Viscosity subsolution]
    We say a hyperbolic convex function $v \in C^0(\mathbb{H}^n)$ is a viscosity subsolution to equation
    \begin{align}
        F(h_{ij}-h\delta_{ij}) = -\frac{h}{\binom{n}{k}^{\frac{1}{k}}}, \label{viscosity subsolution equation}
    \end{align}
    if for any point $y \in \mathbb{H}^n$, and every $\varphi \in C^2$ defined in a neighborhood of $y$, such that $v-\varphi$ has a local maximum at $y$ and
    \[
    v(y) = \varphi(y),
    \]
    we have
    \begin{align*}
        F(\varphi_{ij}-\varphi\delta_{ij}) \geqslant -\frac{\varphi}{\binom{n}{k}^{\frac{1}{k}}}
    \end{align*}
    at $y$.
\end{definition}

\begin{lemma}\label{smooth function g}
    For $0 < a \leqslant \frac{n-1}{n}$, the function 
    \begin{align}\label{g=}
        g(\rho) = - \left( \cosh \rho \right)^{-a}
    \end{align}
    satisfies the above conditions, i.e. $g$ is admissible and is a subsolution for the self-shrinker equation. Moreover, $g$ is hyperbolic convex for any $a>-1$ and $g$ is a strict subsolution for $0 < a < \frac{n-1}{n}$.
\end{lemma}

\begin{lemma}[Viscosity comparison principle]\label{Viscosity comparison principle}
    Suppose $w$ is a viscosity subsolution of \eqref{viscosity subsolution equation}.
    Suppose $v \in \mathcal{S}$ is a $C^2$ supersolution of \eqref{viscosity subsolution equation}.
    If $w \leqslant v$ on $\mathbb{S}^{n-1}$ in cone topology sense, then $w \leqslant v$ in $\mathbb{H}^{n}$.
\end{lemma}

\begin{proof}
    We first claim that the viscosity comparison principle holds for any geodesic ball $B_R$. Suppose that $w$ and $v$ are viscosity subsolution and admissible $C^2$ supersolution in $B_R$ respectively, with
    \begin{align}
        w \leqslant v \quad \text{on } \partial B_R. 
    \end{align}
    We will prove that
    \begin{align}
        w \leqslant v \quad \text{in } B_R. \label{BR leq}
    \end{align}
    Let $\rho$ be the geodesic distance function from the origin. Define
    \[
    g(\rho) = - \left( \cosh \rho \right)^{-a}
    \]
    as in Lemma \ref{smooth function g}, then $g$ is hyperbolic convex, and $g$ is bounded. By Lemma \ref{smooth function g}, we can see that    
    \[
    G(g)=F(g) + \frac{g}{\binom{n}{k}^{\frac{1}{k}}}\geqslant c_0 \quad \text{in } B_R,
    \]
    for some positive constant $c_0$. 
 
    Consider 
    \[
    w_\eta = w + \eta g.
    \]
    We claim that $w_\eta$ is a viscosity subsolution to equation
    \begin{align}\label{Modified subsolution}
        G(h) = \eta c_0.
    \end{align}
    For any $x_0 \in B_R$, suppose  $\varphi_{\eta} \in C^2$ satisfies 
    \[
    \varphi_{\eta}(x_0) = w_{\eta}(x_0), \quad \varphi_{\eta}(z) \geqslant w_{\eta}(z), \quad \mathbb\forall z \in B_R.
    \]
    Let $\varphi = \varphi_{\eta} - \eta g$, then
    \[
    \varphi(x_0) = w(x_0), \quad \varphi(z) \geqslant w(z), \quad \mathbb\forall z \in B_R.
    \]
    Since $w$ is a viscosity subsolution in $B_R$, we have
    \[
    G(\varphi)(x_0) \geqslant 0.
    \]
    By concavity we know at $x_0$
    \begin{align*}
        G(\varphi_{\eta}) \geqslant G(\varphi) + G(\eta g) \geqslant \eta c_0.
    \end{align*}
    By the hyperbolic convexity of $w$, we know that in $B_1^\xi$ model, $w \sqrt{1-|\xi|^2}$ is convex, which implies 
    $(w +\eta g)\sqrt{1-|\xi|^2}$ is also convex. Thus, $w_{\eta}$ is a viscosity subsolution.     
    
    
    Now suppose \eqref{BR leq} fails. We may assume  that there exists $\eta>0$ such that 
    \[
    \max_{\overline{B}_R} \left( w_{\eta}-v \right) > \max_{\partial B_R} \left( w_{\eta}-v \right).
    \]
    Suppose at $z_0 \in B_R$, we have
    \[
    w_{\eta}(z_0)-v(z_0) = M = \max_{\overline{B}_R} \left( w_{\eta}-v \right).
    \]
    Thus we have
    \[
    v(z_0) + M = w_{\eta}(z_0), \quad v + M \geqslant w_{\eta} \quad \text{in } \overline{B}_R.
    \]
    Therefore since $w_{\eta}$ is a viscosity subsolution to \eqref{Modified subsolution}, at $z_0$, we have
    \[
    G(v+M) \geqslant \eta c_0,
    \]
    which shows that at $z_0$
    \begin{align*}
        G(v) = G((v+M) + (-M)) \geqslant G(v+M) + G(-M) \geqslant \eta c_0.
    \end{align*}
    This contradicts $G(v) \leqslant 0$ since $v$ is a supersolution.

    Now we can prove the viscosity comparison principle. Suppose not, then there exists a point $y_0 \in \mathbb{H}^n$, such that $w(y_0) - v(y_0) = 2\varepsilon > 0$. By the same argument as in Lemma \ref{Comparison principle}, we may still assume that $(w-v) \leqslant \varepsilon$ outside $B_R$. Then in $\overline{B}_R$, we have 
    \begin{align*}
        w - \varepsilon \leqslant v \quad \text{on } \partial B_R,
    \end{align*}
    and $w - \varepsilon$ is a viscosity subsolution in $\overline{B}_R$. Using the viscosity comparison principle for geodesic ball $B_R$, we have
    \[
    w - \varepsilon \leqslant v \quad \text{in } B_R.
    \]
    Contradiction. The viscosity comparison principle is proved.
\end{proof}

\end{document}
